\section{\csms}\label{csmodes}

The Channel Strip section has different modes, which can be
selected with the six buttons in the VPOT ASSIGN section.

To switch the modes, press the corresponding VPOT ASSIGN button. When
a channel strip element is used while the VPOT ASSIGN button is
pressed, the original mode will be restored after releasing the VPOT
ASSIGN button. E.g. when the current mode is the \send mode, you can
press and hold the \pan button, use the \select buttons of the channel
to select a different channel and release the \pan button to edit the
send levels of this channel.

The LED of the VPOT ASSIGN button flashes when the mode corresponding
to this button is activated. If it makes sense to select a mode the LED
is turned on. For example when only one track is selected and it contains an
FX, the LED of the \plug button is turned on. The Channel Strip-, Plug-,
Action- and Pan-Mode come with a separate editor. The editor can be opened
by pressing \alt and the corresponding VPOT ASSIGN button.

The Plug-, Action- and Pan-Mode also come with mode specific
options\footnote{The Action- and Pan-Mode options are the same.} that
can be changed with the \vpots while pressing the \option button
(\shift + \option displays additional options).  While you press the
\option button you will see that the MCU display switches. It will
show you up to four properties. In the upper row you will see the
behavior you can change, in the lower row you will see the current
behavior. You can select the behavior with the \vpots. Pressing the
\vpot on the left below an option selects the previous option,
pressing the \vpot on the right below selects the next
option. Alternatively, you can just turn the \vpots. These options
work on a global, not a project-dependent basis.

\noindent
\textbf{Commands.}\label{commands} In addition to the direct functions of
the buttons and \vpots, several channel-strip modes support short
\emph{commands} that are triggered by pressing a \vpot while holding the
\control modifier. The set of commands depends on the mode: in the
\hyperref[panmode]{Pan-Mode} they act on the track of the pressed channel,
in the \hyperref[plugmode]{FX-Mode} and the
\hyperref[channelstrip]{Channel Strip-Mode} they act on a plugin of the
track (opening the floating window or the FX chain, moving the FX up or down
in the chain, removing the FX, ...). While you hold the \control modifier,
the display shows a short label on every \vpot that carries a command, so
you always see which command is available on that \vpot. The command set of
each mode is listed in the section of that mode.

\noindent
\textbf{Modifier keys.} Across the channel-strip modes the four modifier
buttons follow the same convention. \shift is reserved for \emph{addressing}:
it selects the second 8 of a group (e.g.~the \vpots 9\,--\,16, or the FXs,
presets or pages 9\,--\,16) and is never used to trigger a command. \option
switches the display to the mode \emph{options} (with \shift \option the 2nd
options). \alt opens the mode's \emph{editor} when pressed together with the
corresponding VPOT ASSIGN button. And \control triggers the \emph{commands}
described above when pressed together with a \vpot.

\noindent
\textbf{Multiple units.} When more than one hardware unit is
configured (see section~\ref{units}), all channel-strip modes span the
full channel count: each additional unit simply adds eight more
channel strips to the right, so a 16-channel surface shows
tracks~1\,--\,16 at once. The \bankdu and \channeldu buttons move the
whole channel window across all units together. The FX-Mode uses the same
multi-unit channel surface; parameter banks and pages are handled within
the FX-Mode as described in section~\ref{plugmode}.


\pikText{panmode}
\pikText{channelstrip}
\pikText{actionmode}
\pikText{sendmode}
\pikText{plugmode}
